% =====================================================================
%  GMN2 §1  叙述式目录（A narrative table of contents）
%  原文：wkb-paper.tex 325–1031 行；PDF（印刷）p.4–14；编号公式 (1.1)–(1.9) 共 9 条；脚注 11 条
%  本块全部公式均在块内，故 \eqref 可用；跨块引用暂写原文编号并加 % TODO
% =====================================================================

\section{叙述式目录}
\label{sec:IntroductionAndSummary}

超对称规范理论一直是理论物理与数学两方面丰沛的美妙惊喜之源。其中格外丰富的一类，是具有 $d=4$、$\CN=2$ 超对称的理论（下称 $\CN=2$ 超对称）。在这一背景下，希尔伯特空间中一个重要的、被单独挑出的子空间是 BPS 态的空间。这些态处于超对称代数的「小」表示或「刚性」表示之中；正是这种刚性使它们便于分析，而这又成了这些理论已知精确结果的基础。然而，尽管 $\CN=2$ 的 BPS 表示是刚性的，BPS 子空间仍然同时非平凡地依赖于紫外参数与该理论量子真空的选取——这一现象称为穿墙。

尽管 BPS 子空间如此基本，给定一个任意的 $\CN=2$ 理论及其真空选取，却并没有把它算出来的算法。事实上，BPS 谱只在少数例子中被显式知道，而那里都要用到专门的\emph{特设}技巧。本文的主要结果是一个新算法，用以确定某一无穷族 $\CN=2$ 理论的 BPS 谱。我们的方法所适用的理论在第 3 节中描述% TODO 合稿改 \ref{sec:physics}
，新算法在本文第 11 节中描述% TODO 合稿改 \ref{sec:UniversalStokesMatrix}
。这个结果之所以值得期待，是因为有迹象表明，该算法的推广将适用于远为广泛的一类 $\CN=2$ 理论。

我们的主要结果所适用的那类理论，是各节点处规范群为 $SU(2)$ 因子的线性箭图规范理论。它们是更大的一类特殊 $\CN=2$ 规范理论的一部分，第 3 节对此有详细描述。这个更大的理论类我们称为 $\CS$（取「六」之意），它来自 M5-膜在带穿刺点的黎曼面 $C$ 上的紧化。\footnote{更准确地说，我们考虑的是由引力脱耦而得到的低能 $(2,0)$ 超共形场论；随后我们把该理论在 $C$ 上紧化，并对 $d=6$ 的 $(2,0)$ 超代数作部分扭转，以保持 $\CN=2$ 超对称。} 超共形 $(2,0)$ 理论有 A-D-E 分类，因此 $\CS$ 中的元素可由一个单连通的「规范群」、一个黎曼面 $C$、以及用「缺陷算符」对 $C$ 的穿刺点所作的一种装饰来标定。$\CS$ 类中的理论有一个重要性质：它们与 Hitchin 系统关系密切。这一关系是本文的绝对核心，它在再进一步紧化到圆上的三维时显现出来。低能下，三维有效理论是一个靶空间为 $\CM$ 的 $d=3$、$\CN=4$ sigma 模型。作为黎曼流形，这个靶空间可以等同于某个 Hitchin 系统之解构成的模空间。要说明这一点，关键的观察是：不必先在 $C$ 上、再在 $S^1$ 上紧化；我们可以——借助场论版本的「Fubini 定理」——先在 $S^1$ 上、再在 $C$ 上紧化，从而构造出同一个三维有效理论。先对 $S^1$ 紧化给出五维超对称 Yang–Mills 理论；随后把（扭过的）$d=5$ 超 Yang–Mills 理论在 $C$ 上紧化，所得的 BPS 方程正是众所周知的 Hitchin 方程。特别地，若从 $K$ 个 M5-膜出发（即六维中的超共形 $u(K)$ $(2,0)$ 理论），则 Hitchin 方程就是下面的 \eqref{eq:hitchin-1}–\eqref{eq:hitchin-3}：

\begin{align}
F + R^2 [\varphi , \bar\varphi] & = 0, \label{eq:hitchin-1}\\
\partial_{\bar z} \varphi  + [ A_{\bar z}  ,\varphi] & = 0, \label{eq:hitchin-2}\\
\partial_z \bar\varphi  + [ A_z,\bar\varphi]  & = 0, \label{eq:hitchin-3}
\end{align}

其中 $R$ 是圆的半径，$F$ 是 $C$ 上 $u(K)$ 规范场 $A$ 的场强，$\varphi$ 是取值于伴随丛的 $1$-形式的 $(1,0)$ 部分；$z$ 是 $C$ 上的局部全纯坐标。

若不指定 $(A,\varphi)$ 在 $C$ 的穿刺点处的边界条件，对 Hitchin 系统的描述就是不完整的。在这些穿刺点处，场 $(A,\varphi)$ 带有奇性。物理上，这些奇性编码了六维超共形理论中那些多少有些神秘的「缺陷算符」（而实际上，缺陷算符正是由 $(A,\varphi)$ 被指定的奇性来定义的）。要考虑的最简单的算符——也是我们关注的那些——来自多重 $u(K)$ M5-膜理论与「横向」单重缠绕 M5-膜在 $C$ 的穿刺点处的相交。所谓横向，指的是下面这个意思。一般而言，曲线 $C$ 作为一条全纯曲线嵌在某个超凯勒流形 $Q$ 中；引力脱耦极限允许我们把 $Q$ 换成 $T^*C$ 零截面的一个邻域。横向的五膜填满 $\CN=2$ 理论的四维时空 $\IR^{1,3}$，并沿着投影 $T^*C \to C$ 的纤维延伸。在第 3 节中我们说明如何把这一物理图像翻译成对 $(A,\varphi)$ 之奇性的条件。这些奇性在第 3.2.4 节中有详细描述：正则奇性的情形见 (3.74)、(3.75)、(3.76)；% TODO 跨块引用
不规则奇性的情形见第 3.2.6 节的方程 (3.115)、(3.116) 以及第 9.3 节。

Hitchin 系统在全文各处都起核心作用，本文第 4 节% TODO \ref{sec:hitchin}
概述了我们所需关于 Hitchin 系统的基本事实。偏数学的读者可以跳过第 3 节，直接读第 4 节的简要总结——只不过那样一来，寻找 BPS 态的规则可能会显得有些缺乏动机。

$\CS$ 类中理论的一批格外重要的例子，由 Witten 用 NS5-膜与 D4-膜阵列对 $\CN=2$ 理论所作的几何构造给出\cite{Witten:1997sc}。这些构造常常用图 5% TODO 图号待合稿
这样的图来概括。第 3 节的大部分内容仅是对 Witten 构造的回顾，熟悉文献 \cite{Klemm:1996bj,Witten:1997sc} 的读者可以安全跳过。不过我们要指出，第 3.2.4、3.2.5、3.2.6 节% TODO 跨块引用
中含有一些新的要点，讲的是如何用物理图像来描述 $(A,\varphi)$ 的边界条件。此外，第 3.2.8 与 3.2.9 节% TODO 跨块引用
分别含有关于 Hitchin 模空间之间的同构、以及关于味对称的一些新评论。特别地，第 3.2.8 节的同构应当具有一定的数学意义；Philip Boalch 最近独立地注意到了类似的一类同构\cite{irkm}。

$C$ 上 Hitchin 系统的几何，优美地编码了 $\CS$ 中\emph{四维} $\CN=2$ 理论低能有效 Seiberg–Witten 理论的关键数据。首先，Seiberg–Witten 曲线——即 $C$ 的一个分歧覆盖 $\Sigma$——恰好以 Hitchin 系统谱曲线的身份漂亮地呈现出来。于是 $\Sigma \subset T^*C$，而且覆盖的各叶可以由 $\varphi$ 的本征值来标记。相当一般地，余切丛 $T^*C$ 上典范地带有某个辛形式，而这个辛形式又被一个一形式典范地平凡化；把这个一形式限制到 $\Sigma$ 上，就得到 Seiberg–Witten 微分，记作 $\lambda$。

既然本文的主题是 BPS 谱，我们理应弄懂如何在 Hitchin 系统的语境中描述这个谱。既然这些理论源自缠绕在 $C$ 上的 M5-膜，人们可以系统地用六维理论中的弦来理解 BPS 态，而这些弦又来自终止于 M5-膜上的开 M2-膜\cite{Strominger:1996ac}。把对 BPS 态的这种描述翻译成 Hitchin 系统的语言，见第 3.1.1 与 3.1.3 节% TODO 跨块引用
。在 $K=2$ 的情形，我们回到 Klemm 等人一个众所周知的构造\cite{Klemm:1996bj}：BPS 态与 $C$ 上使弦的局部张力极小的曲线相关联。用数学的话说，在曲线上一点 $z_0$ 附近，我们选取覆盖的一个单值分支，并用 $w = \int^z_{z_0} \lambda$ 定义局部坐标；于是该曲线在 $w$-平面中是一条直线。BPS（半）超多重态与那些起止于覆盖 $\Sigma \to C$ 的分支点上的曲线相关联，而 BPS 矢量多重态与那些闭合的曲线相关联。出于下面将要解释的理由，这两类曲线我们都称为\emph{有限 WKB 曲线}。插图见图 3 与图 4。% TODO 图号待合稿
相应 BPS 态的中心荷是 $Z_\gamma = \pi^{-1} \oint_{\gamma} \lambda$。下面一个要点是：这个中心荷的相位，就是 $w$ 平面中该直线与 ${\rm Re}(w)$ 轴之间的夹角 $\vartheta$。当 $K>2$ 时，对 BPS 态的类似描述要更复杂一些，并要用到 $C$ 上的弦网。

正如常有的情形，我们的主要结果其实并不是这项工作最初的目标。最初的动机来自最近一个关于超凯勒度量的构造\cite{Gaiotto:2008cd}，而那个构造本身又是为 Kontsevich–Soibelman 穿墙公式\cite{ks1}寻找物理根基的努力所促成的。为使本导读余下的部分可读，我们必须先在此回顾 \cite{Gaiotto:2008cd} 中最本质的几点；更技术性的总结见本文第 2 节% TODO \ref{sec:review}
。

\cite{Gaiotto:2008cd} 的出发点是：把一个一般的 $d=4$、$\CN=2$ 理论放在半径为 $R$ 的圆上紧化。如上面提到的，低能下该理论是一个三维 sigma 模型，其靶空间 $\CM$ 必须带有超凯勒度量。空间 $\CM$ 有一个纤维化 $\CM \to \CB$，其中 $\CB$ 是四维理论的真空模空间，而一般的纤维是维数为 $\dim \CB$ 的实环面。当 $R\to\infty$ 时，$\CM$ 上的度量指数式地接近于一个简单而显式的度量；该度量是超凯勒的，但在实余维二处有奇性。这个度量称为「半平坦度量」，记作 $g^{\sf}$。它很容易由 Seiberg–Witten 有效拉氏量沿 $S^1$ 的朴素量纲约化得到。$R$ 取有限值时的量子修正把 $g^{\sf}$ 抹平。此外，这些量子修正只依赖于 BPS 态的谱；而这个谱本身又是（四维）真空的函数，所以度量的光滑性蕴含着一个穿墙公式。在 \cite{Gaiotto:2008cd} 中证明了这正是 Kontsevich–Soibelman 穿墙公式（KSWCF）。

从纯粹数学的观点看，可以把 \cite{Gaiotto:2008cd} 视为由下面三件数据出发构造超凯勒度量：

\begin{itemize}

\item D1：格的局部系统 $\hat \Gamma\to \CB$，其纤维上带一个整值反对称形式 $\langle,\rangle$（可以退化）。令 $\Gamma$ 为用 $\hat \Gamma$ 的根除所得的辛格局部系统，则要求 $\CB$ 的实维数等于 $\Gamma$ 的秩。

\item D2：中心荷函数 $Z \in {\rm Hom}(\hat \Gamma, \IC)$，满足 $\langle dZ, dZ \rangle = 0$，其中 $d$ 是沿 $\CB$ 的微分。

\item D3：分段常值函数 $\Omega: \hat \Gamma \to \IZ$，满足 KSWCF。

\end{itemize}

当一个物理的 $\CN=2$ 理论提供数据 D1、D2、D3 时，该物理问题的模空间 $\CM$ 可以等同于 $\CB$ 上纤维化 $\Gamma^* \otimes \IR/2\pi \IZ$ 的全空间。于是我们预期：仅给出 D1、D2、D3，就能在 $\CM=\Gamma^* \otimes \IR/2\pi \IZ$ 上构造出超凯勒度量——而事实证明正是如此。实际上，我们构造的是 $\CM$ 上一\emph{族}超凯勒度量；族的参数下面会说明。

构造 $\CM$ 上度量的一个本质部分，涉及超凯勒度量的扭量描述。利用 $\hat \Gamma^* \otimes \IR/2\pi \IZ$ 具有环面纤维化这一事实，可以把扭量空间上适当全纯数据的构造，归结为构造某个映射 $\CX^\RH: \CM \times \IC^\times \to \hat \Gamma^* \otimes \IC^\times$\footnote{在 \cite{Gaiotto:2008cd} 中这个映射就叫做 $\CX$。}，其中定义域中的因子 $\IC^\times$ 是扭量球的一部分。把 $\CX^\RH$ 与 $\gamma\in \hat \Gamma$ 缩并，记作 $\CX^\RH_{\gamma}: \CM \to \IC^\times$，这会非常方便。我们把这些函数称为\emph{达布坐标}\footnote{这里有一处我们引以为憾的术语滥用。一旦选取 $\Gamma$ 的一组与某个拉格朗日分解相容的基 $\{\gamma_i\}$，并把它提升到 $\hat \Gamma$，则 $\log \CX_{\gamma_i}^\RH$ 才真正给出全纯辛流形 $\CM$ 上的一组达布坐标。}。注意 $\CX^\RH_{\gamma}\CX^\RH_{\gamma'} = \CX^\RH_{\gamma+\gamma'}$。映射 $\CX^\RH$ 受一列定义性质的约束。详细性质的完整清单在第 2 节中回顾，但有三个关键性质必须在此陈述：

\begin{itemize}

\item P1：第一，泊松结构由方程 (\ref{eq:poisson}) 定义：
\begin{equation}
\{ \CX^\RH_{\gamma}, \CX^\RH_{\gamma'} \} = \langle \gamma,\gamma' \rangle \CX^\RH_{\gamma+\gamma'}. \label{eq:poisson}
\end{equation}

\item P2：第二，$\CX^\RH_{\gamma}$ 渐近于与半平坦度量 $g^{\sf}$ 相伴的类似函数 $\CX_{\gamma}^{\sf}$。设 $\zeta \in \IC^\times$ 是扭量球上的一个元素，则半平坦达布坐标可以非常显式地写成
\begin{equation}\label{eq:P2-sf}
\CX_{\gamma}^{\sf} := \exp \left( \pi R\zeta^{-1} Z_\gamma + i \theta_\gamma + \pi R \zeta \bar Z_\gamma \right),
\end{equation}
其中 $\theta_\gamma: \hat \Gamma^*\otimes \IR/2\pi \IZ \to \IR/2\pi \IZ$ 由缩并典范地定义。我们要求 $\CX_{\gamma}^\RH \sim \CX_{\gamma}^{\rm sf}$，\emph{既}在 $\zeta\to 0, \infty$ 时、\emph{也}在 $R\to \infty$ 时成立。

\item P3：第三，$\CX^\RH$ 作为 $\zeta$ 的函数，其解析结构受到如下约束。称射线 $\ell_{\gamma,u} := \{ \zeta: Z_\gamma(u)/\zeta \in \IR_-\}$ 为\emph{BPS 射线}。于是，当 $\zeta$ 穿过一条 BPS 射线 $\ell_{\gamma_0,u}$ 时，$\CX^\RH_{\gamma}$ 会有一个由泊松变换 $\CK_{\gamma_0}^{\Omega(\gamma_0;u)}$ 给出的间断，其中\footnote{其实应当考虑 $\gamma_0$ 的所有倍数，因此正确的变换是 $\prod_{\gamma_0'' \parallel \gamma_0} \CK_{\gamma_0''}^{\Omega(\gamma_0'';u)}$。在我们研究的例子里，只有单个电荷对该间断有贡献。}
\begin{equation}
\CK_{\gamma_0} : \CX^\RH_\gamma \to \CX^\RH_{\gamma} (1 \pm \CX^\RH_{\gamma_0} )^{\langle \gamma, \gamma_0 \rangle}. \label{eq:KS-Transfm}
\end{equation}
变换 $\CK_{\gamma_0}$ 将被称为 \emph{KS 变换}。更精确的方程（尤其是 $\pm$ 号的选取）见下面的 (\ref{eq:KS-Transfm})。此外，$\CX_\gamma^\RH$ 作为 $\zeta$ 的函数，必须在带 $\Omega(\gamma;u) \not=0$ 的那些 BPS 射线 $\ell_{\gamma,u}$ 之集合的补集上\emph{全纯}（不带任何奇性）。
\end{itemize}

如 \cite{Gaiotto:2008cd} 中所解释，由函数 $\CX^\RH_{\gamma}$ 可以恢复 $\CM$ 上的超凯勒度量。所得度量把 $g^{\sf}$ 的实余维二奇性抹平（但某些实余维四的奇性可能仍在）。

回到物理的观点，函数 $\CX^\RH_{\gamma}$ 有很好的解释：它们既可以在三维拓扑场论中表示成线算符的期望值，也可以看作某个手征环的元素。我们希望另文描述 $\CX^\RH_\gamma$ 的这些方面（附录 A 的注 3 中我们也会再简短提及这一点）。

上面我们说得到了一族超凯勒度量。要理解这一点，首先注意 $\hat \Gamma$ 的根，物理上就是味荷格。我们有
\begin{equation}
0 \rightarrow \Gamma_{\rm flavor} \rightarrow \hat \Gamma \rightarrow \Gamma \rightarrow 0, \label{eq:flavor-seq}
\end{equation}
而辛格 $\Gamma$ 是电磁规范荷格。流形 $\hat \Gamma^*\otimes \IR/2\pi \IZ$ 被 $\CM$ 的许多拷贝叶层化。此外，对 $\gamma \in \Gamma_{\rm flavor}$，$\CX^\RH_{\gamma}$ 取 (\ref{eq:P2-sf}) 的精确形式，其中 $Z_\gamma$ 编码超多重态质量、$\theta_\gamma$ 编码味 Wilson 线。这些参数连同 $R$ 一起，参数化了 $\CM$ 上那族超凯勒度量。

关于 \cite{Gaiotto:2008cd} 我们必须回顾的最后一点是 $\CX^\RH$ 的显式构造。做法是：利用性质 P2 与 P3，把 $\CX^\RH$ 刻画成某个黎曼–希尔伯特问题的解，而该问题又等价于一个积分方程；再通过一系列操作把它解出来。顺便说一下，这个积分方程其实是热力学贝特拟设（TBA）的一个版本\footnote{四维超 Yang–Mills 理论与 TBA 之间的另一个关系最近由 Nekrasov 与 Shatashvili 讨论过\cite{Nekrasov:2009ui}。}。\emph{只要}半径 $R$ 很大，这个 TBA 方程就可以用迭代求解。结果是多重积分给出的显式级数展开，当 $\CX_{\gamma}^{\sf}$ 很小时其被积函数也很小。于是我们得到超凯勒度量的显式构造。必须强调，$R$ 应当很大，才能保证用迭代解积分方程是合理的。（这里用到的一个重要事实是：当 $R\to \infty$ 时 $\CX_\gamma^{\sf}$ 在 $\ell_{\gamma,u}$ 上是指数式小的。）

回顾完 \cite{Gaiotto:2008cd} 的主要特征，我们终于可以回到撰写本文的原初目标：给出 $\CX^\RH_{\gamma}$ 的另一种构造，它不依赖那个积分方程，且对一切 $R$ 成立。本文确实将给出与 $SU(2)$ Hitchin 系统相伴的 $\CS$ 类理论中函数 $\CX^\RH_{\gamma}$ 的另一种构造。我们的定义对\emph{一切 $R$ 值}都合理；而且这个新构造还便于对关键定义性质 P1、P2 以及 P3 的（一部分）作优雅的几何验证\footnote{熟悉热力学贝特拟设的读者可以把我们的结果看作对 \cite{Dorey:1999uk} 的一个宽泛推广：在那里，TBA 方程的解（在「共形极限」下）基本上是由某个全纯联络的单值数据构造出来的。}。

在解释这个新构造之前，我们必须一开始就坦白：有一个困难将仍然悬而未决。关于达布坐标在小 $R$ 时的行为，本文与 \cite{Gaiotto:2008cd} 之间有些不一致。在 \cite{Gaiotto:2008cd} 中我们设想 TBA 方程对一切 $R$ 都有正则解，那样就会给出一个在 $\zeta$-平面上没有极点的 $\CX^\RH$\footnote{更准确地说，$\CX^\RH$ 在带 $\Omega\not=0$ 的 BPS 射线之集合的补集上是全纯的。}。本文的结果提示真相也许更复杂：我们确实找到了 $\CX^\RH$ 的一个自然候选，而且它确实对一切 $R$ 都有定义；但当 $R$ 小时，我们无法证明它无极点。另一方面，也很难设想存在两个\ti{不同}的 $\CX^\RH$、一个有极点一个没有的情形。于是我们看到两个合理的选项。其一，本文的 $\CX^\RH_\gamma$ 其实没有极点；那样它们对一切 $R$ 都能与 \cite{Gaiotto:2008cd} 积分方程的解等同起来。其二，本文的 $\CX^\RH_\gamma$ 确有极点；那么严格说来，它们在小 $R$ 时就不是 \cite{Gaiotto:2008cd} 那个积分方程的解。这种「额外极点」随参数变化而出现的现象，在 TBA 的文献中早已为人熟知（例如见 \cite{Dorey:1996re}），并且在那个语境下是可以处理的。显然，这些问题值得进一步关注！

现在让我们综述 $\CX^\RH_{\gamma}$ 的新构造。关键的想法始于这样一个事实：Hitchin 模空间 $\CM$ \emph{同时}是绕 $C$ 的穿刺点具有固定单值\footnote{由于 Hitchin 系统可以同时带有正则与不规则奇性，这里的「单值数据」包含 Stokes 数据。}的平坦联络的模空间。确实，给定 Hitchin 方程的一个解，我们可以构造一个非酉联络（即下面的 \eqref{eq:def-CA}）
\begin{equation}
\CA : = \frac{R}{\zeta} \varphi + A + R \zeta \bar \varphi, \label{eq:def-CA}
\end{equation}
而 Hitchin 方程蕴涵它对任意 $\zeta \in \IC^\times$ 都是 $C$ 上的平坦联络；反之，这样一个联络对一切 $\zeta$ 的平坦性也蕴涵 $(A,\varphi)$ 解 Hitchin 方程。下一个观察是：Fock 与 Goncharov 利用 $C$ 的一个三角剖分的数据，构造了平坦联络模空间上一组漂亮的坐标\cite{MR2233852}。（这组有用的坐标只在 $C$ 至少有一个穿刺点时才存在；因此即使物理理论在 $C$ 没有穿刺点时也有意义，我们也不打算把 $\CX_\gamma$ 的构造推广到那种情形之外。）我们将用 Fock–Goncharov 坐标来构造函数 $\CX_{\gamma}$。我们的纲领大致如下：第一步，给定一个角\footnote{$\vartheta$ 的周期性可以是 $2\pi$ 的整数倍，它甚至可以取值于万有覆叠 $\IR$。} $\vartheta$，我们定义一个特殊的三角剖分，称为\emph{WKB 三角剖分}；第二步，把 Fock–Goncharov 构造应用于该三角剖分，得到一组函数 $\CX_{\gamma}^{\vartheta}: \CM \times \IC^\times \to \IC^\times$；接着令 $\vartheta = {\rm arg}(\zeta)$，用 $\CX_{\gamma}^{\vartheta}$ 反过来构造 $\CX^\RH_{\gamma}$；第三步，证明所得函数满足第 2 节中概述的那些定义性质（尤其是 P1、P2、P3）。下面我们稍详细地勾勒一下这个纲领是如何实现的。

在第 5 节% TODO \ref{sec:fg-coordinates}
中我们回顾 Fock 与 Goncharov 的构造\cite{MR2233852}。我们在两处偏离了他们的讨论。一个相当次要的差别是：我们更愿意使用\emph{装饰三角剖分}（decorated triangulation）。在正则奇点（我们主要关注的情形）下，这就是顶点取在 Hitchin 系统奇点 $\CP_i$ 处的理想三角剖分，但我们在每个点 $\CP_i$ 处多加一件数据。具体说，我们考虑解 $(d+\CA)s=0$ 的平坦截面 $s$，而 $\CP_i$ 处的装饰就是在其邻域内选一个只定到标度因子的平坦截面 $s_i$。这样的平坦截面必然是绕 $\CP_i$ 的单值的本征向量；等价地说，装饰就是在该单值的两条本征线中选定一条。（对不规则奇点的类似「装饰三角剖分」概念，在第 8 节% TODO \ref{sec:irregular}
中解释。）与 Fock 和 Goncharov 的工作相比第二个、也是更重要的偏离是：BPS 谱中矢量多重态的存在，迫使我们把三角剖分的概念扩充，以容纳更精细的对象，我们称之为「极限三角剖分」（limit triangulation）。这些在第 5.9 节% TODO
中描述。Fock–Goncharov 构造的核心，是用（平行移动后的）平坦截面 $s_i$ 之间的「交叠」来描述平坦联络 $\CA$ 的单值。在附录 A% TODO \ref{app:MonodromyX}
中我们说明如何做到这一点；这一程序自然导致方程 (5.2) 中那个关键定义。% TODO 跨块引用

$\CX^\RH_\gamma$ 的性质 P1 会正确地显现出来，这一点在这个阶段已经可以看得很清楚。Hitchin 模空间带有一个自然的全纯辛形式，即 (4.10) 中给出的% TODO
\begin{equation}
\varpi_\zeta = \frac{1}{2}\int_{C} {\Tr} \delta \CA \wedge \delta \CA. \label{eq:symplectic}
\end{equation}
Fock–Goncharov 坐标相应的泊松括号于是取简单形式 (5.9)。% TODO 跨块引用
为自足起见，我们在第 5.4–5.5 节% TODO
中用附录 B% TODO \ref{app:poorman}
的结果给出这些泊松括号的一个初等推导。

坐标 $\CX^\RH_{\gamma}$ 的第二个关键定义性质（上面的 P2）是它们在 $\zeta \to 0,\infty$ 与 $R \to \infty$ 时的渐近行为。正是这个性质促成了我们对 WKB 三角剖分的定义。如第 6 节% TODO \ref{sec:WKB-triangulation}
所述，我们定义相位为 $\vartheta$ 的 WKB 曲线满足 $\langle \lambda, \partial_t \rangle = e^{i \vartheta}$。当然，我们在上面讨论 BPS 态时已经遇到过这个条件！它等价于说：在局部坐标 $w = \int^z_{z_0} \lambda$（其中 $z_0$ 是曲线上一点）中，该曲线是一条平行于 $e^{i \vartheta}$ 的直线。这些 WKB 曲线倾向于被奇点「俘获」（如第 6.2 节% TODO \ref{sec:wkb-local}
的局部分析所示），因此一般的 WKB 曲线起止于奇点。WKB 三角剖分于是通过选取适当一个有限的、「拓扑上互不相同」的一般 WKB 曲线集来定义，所用程序在第 6.3 与 6.4 节% TODO
中解释。此外还有一个典范的装饰：选取沿一条坠入奇点的 WKB 曲线指数式变小的那个平坦截面。这一装饰三角剖分的选取，是由对平坦截面方程 $(d+\CA)s=0$ 作 WKB 分析（小参数为 $\zeta$、$1/\zeta$ 或 $1/R$）所促成的。回忆一下：在 WKB 近似中，指数式小的波函数可以可靠地算出，而指数式大的波函数则因为可以加上一个未知的指数式小分量而有歧义。正因如此，在计算从 $C$ 上不同位置 $\CP_i$ 平移来的平坦截面 $s_i$ 的「交叠」时必须小心：我们必须沿着 WKB 曲线来平移这些截面。总之，给定一个角 $\vartheta$ 与一个真空 $u\in \CB$——或者更好，给定一个 Seiberg–Witten 微分 $\lambda$——我们就有一个典范确定的装饰三角剖分。它记作 $T_\WKB(\vartheta,u)$ 或 $T_\WKB(\vartheta, \lambda^2)$；后一种记法反映了这样一个事实：边是无向的，因而只依赖于二次微分 $\lambda^2$。

现在转向第三个关键性质 P3：我们必须考虑，在固定 $\lambda$ 的情况下改变 $\vartheta$ 时，不同的 WKB 三角剖分之间如何相互关联。相当一般地，不同的装饰三角剖分（不一定是 WKB 型的）之间可以通过一系列初等变换互相转换，我们把这些变换称为 \emph{flip}、\emph{juggle} 与 \emph{pop}。% 译者注：这三个词是 Fock–Goncharov 与本文作者的术语造词，保留原文
我们可以把装饰三角剖分看作某个群胚中的对象，而 flip、juggle、pop 就是生成该群胚中所有其他态射的初等态射。flip 就是在两个三角形拼成的四边形内翻转一条边这一标准变换（如图 16 所示），% TODO 图号
装饰保持不变。pop 则不改动三角剖分，只改变某个指定顶点处那条特殊本征线的选取；当 $\CA$ 的结构群是 $SL(2,\IC)$ 且奇点是单值可对角化的正则奇点时，pop 变换就是互换那两条本征线。最难的变换 juggle 则联系不同的极限三角剖分；详细讨论见第 5.9 节，图示见图 29。% TODO
Fock–Goncharov 理论的一个重要方面是：在 flip 之下坐标经历 \emph{cluster 变换}。这些 cluster 变换原来是 Kontsevich–Soibelman 变换 $\CK_{\gamma_0}$ 的特例。pop 之下的变换是显式已知的，但一般相当繁冗。对我们的主要结果而言重要的一点是：尽管如此，\emph{omnipop}——即同时对所有顶点作 pop 的变换——仍是一个简单且可计算的变换 $\mathbf{S}$，我们称之为\emph{谱生成元}，理由下面会清楚。omnipop 变换 $\mathbf{S}$ 在第 11.1 节% TODO \ref{subsec:SuperPop}
中推导；juggle 之下的变换在第 5.9 节中描述。

在第 5 与第 6 节% TODO
中把全部机器搭好之后，我们终于给出函数 $\CX^{\vartheta}_\gamma:\CM \times \IC^\times \to \IC$ 的关键定义，即方程 (7.3) 与 (7.4)。% TODO 跨块引用
由于该定义是用 Fock–Goncharov 坐标给出的，泊松括号（上面的性质 P1）自然随之而来。此外，我们所选的装饰三角剖分 $T_\WKB(\vartheta, \lambda^2)$ 使得 $\zeta \to 0, \infty$ 与 $R \to \infty$ 的渐近行为可以直截了当地推出——这正是它被设计出来要做的事。现在，相当漂亮的是：当 $\vartheta$ 变化时，WKB 三角剖分恰好发生在 $\vartheta$ 等于某条 BPS 射线 $\ell_{\gamma,u}$ 的倾角时经历 flip、juggle 与 pop。这其实很自然，因为如上面所述，\cite{Klemm:1996bj} 的结果简单说来就是：BPS 态与非一般的\emph{有限} WKB 曲线相关联；回忆一下，这些曲线或是闭合的，或是起止于覆盖 $\Sigma \to C$ 的分支点上的。这些基本态射如图 27 与图 29 所示，% TODO 图号
而方程 (7.33) 与 (7.56) 表明：$\CX^{\vartheta}_{\gamma}$ 中相应的间断恰恰就是与 KS 变换相伴随的那些间断，且 $\Omega$ 的值正确。

在第 11 节% TODO \ref{sec:UniversalStokesMatrix}
中，我们用第 7 节的结果给出计算 $K=2$ 理论 BPS 谱的算法。在复 $\zeta$ 平面中取由角 $[\vartheta, \vartheta + \pi]$ 处两条射线所界定的半平面 $\CH(\vartheta,\vartheta+\pi)$，并考虑让三角剖分 $T_\WKB(\vartheta,\lambda^2)$ 随 $e^{i\vartheta}$ 在该半平面内转至 $-e^{i \vartheta}$ 而演化。值得注意的是，在这一连续演化中，pop 总是发生在特殊的场合\footnote{更技术性地说：pop 发生在一个退化三角形的中心，例如图 19 所示。}，使得 $\CX^{\vartheta}_{\gamma}$ 的相应变换恰好是恒等变换。这一出人意料的事实在第 7.6.3 节% TODO \ref{sec:no-flavor-jump}
中被证明。于是，让 $\vartheta \to \vartheta + \pi$ 演化所得的总变换只涉及一列 flip 与 juggle，而它对达布坐标的作用就是变换 ${\bf S} = \prod \CK_{\gamma}^{\Omega(\gamma;u)}$，其中因子按 $\arg Z_\gamma$ 排序。现在，每个 BPS 态（或其反粒子）在所选半平面 $\CH(\vartheta,\vartheta+\pi)$ 中都有一条 BPS 射线，所以这个乘积恰好捕获谱的一半，另一半只是反粒子。另一方面，初、末两个三角剖分 $T_\WKB(\vartheta, \lambda^2)$ 与 $T_\WKB(\vartheta+\pi,\lambda^2)$ 原来只相差一个 omnipop，而如上面所述，那个变换可以显式算出（第 11 节）。既然乘积分解 ${\bf S} = \prod \CK_{\gamma}^{\Omega(\gamma;u)}$ 是唯一的（在给定 BPS 射线次序的前提下，而该次序又由 $u$ 决定），我们就能从名字起得恰如其分的谱生成元 ${\bf S}$ 读出谱来。值得一问的是：这个算法比 \cite{Klemm:1996bj} 中已给出的、用于计算 $\CS$ 类 $A_1$ 理论 BPS 谱的方案好在哪里？后一方案要求人们预先知道使 BPS 态存在的那些 $\vartheta$ 的临界值；上面算法的要点恰恰是：只需选取一个一般的 $\vartheta$ 值，完全不需要预先知道被占据的 BPS 中心荷的相位。

在第 9 与第 10 节% TODO
中我们算出本形式体系的大量例子。在第 9 节中我们说明线性 $SU(2)$ 箭图理论的各种极限如何囊括所有可能的 Argyres–Douglas（AD）超共形理论。AD 理论最简单的例子就已经为 KSWCF 提供了漂亮的例证：原来这些理论中所有的穿墙恒等式都可由一个基本的五边形恒等式 (9.31) 推出。% TODO 跨块引用
在第 10 节中我们考虑带 $N_f=0,1,2,3,4$ 种基本超多重态味的 $SU(2)$ 规范理论。此时 BPS 谱非常精细，穿墙通常涉及 KS 变换的无穷乘积，推广了 $N_f=0$ 时出现的基本例子，即方程 (10.2)。% TODO 跨块引用
结果是：这些例子中所有的穿墙公式——无论多么复杂——都可以由这个基本恒等式与五边形恒等式的反复使用得到。最精细也最漂亮的谱出现在 $N_f=4$ 的情形。我们定位出一个格外有意思的强耦合区域，其中有限的谱分裂成两个 $N=4$ 的 AD 点，因而可以被非常具体地描述。我们预期这一观察在将来某些研究中会有用。

让我们以回到原初目标来结束这篇导读：定义与 \cite{Gaiotto:2008cd} 的构造相关的那些函数 $\CX^\RH_\gamma$。一个重要的预备结果在第 13 节% TODO \ref{sec:R-To-Infinity}
中给出：我们证明，在 $R\to \infty$ 极限下，只要 $\zeta$ 位于以 $e^{i \vartheta}$ 为中心的那个半平面 $\IH_{\vartheta}$ 内，$\CX_{\gamma}^{\vartheta}$ 确实渐近于半平坦坐标 $\CX_{\gamma}^{\sf}$。证明用到与黎曼面 $C$ 上 sinh-Gordon 方程（及其推广）的一个有趣联系\footnote{顺便说，这个关系最近在 Alday 与 Maldacena 的工作\cite{Alday:2009yn}中也起了有益的作用，而正是那个联系使那两位作者得以很好地利用本文的部分结果。我们在第 9.4.3 节% TODO \ref{sec:pnut}
中再稍多讨论一下这个联系。}。由此我们可以推出：$\CX_{\gamma}^{\vartheta = \arg \zeta}$ 的奇性只有 $\zeta = 0,\infty$ 处的本性奇性，以及穿过那些 $\Omega\not=0$ 的 BPS 射线时的间断。结合第 7 节的其他结果可知：它们满足第 2 节中全套定义性质，因而确实提供了我们所要寻求的那种替代构造。令人恼火的是，如我们已经提到的，这一推理对足够大的 $R$ 成立，但在某个有限的 $R$ 可能失效（第 14 节% TODO \ref{sec:comparison}
有解释）。于是，当 $R$ 减小到零时的行为就成了一个重要的公开问题。

最后我们以若干评论与未来可能的研究方向作结：

\begin{itemize}

\item 首先，也最显然的是：应当把结果推广到更高秩的理论（$K>2$）。自本文第一个预印本版本出现以来，我们在这方面已有一些进展；一般的情形将出现在 \cite{gmn5-to-appear} 中，一些重要的例子将出现在 \cite{gmn6-to-appear} 中。

\item 我们的构造提出了一些诱人的可能联系，通往用 Hall 代数计数 BPS 态的数学（也许还能从那里通往 BPS 态代数）。看起来很可能在这套故事背后应当有一个范畴，某种「黎曼面中一维流形的 Fukaya 范畴」。如我们在第 7.8 节% TODO \ref{sec:bps-comments}
中所评论，可以设想这个范畴其实等价于某个箭图表示的范畴。此外，我们所考虑的 $C$ 的分解的几何，提示了一种在该范畴中实现 Harder–Narasimhan 滤链的方法。我们在第 12 节% TODO \ref{sec:category}
中简要勾勒这些想法。

\item 我们结果的一个推论，是对某些 Higgs 丛模空间上超凯勒度量的新描述。Higgs 丛的模空间在 \cite{geom-lang-n4} 所开创的几何 Langlands 纲领的新进路中扮演重要角色；而本文中这些模空间出现的方式，与它们在 \cite{geom-lang-n4} 中出现的方式确实并非无关。因此自然要问：我们的结果对几何 Langlands 是否有用。

\item Fock 与 Goncharov 的构造，部分动机是想要构造新的无穷维模函子（与 Liouville 理论及其更高秩的「Toda」推广相关联）。我们相信本文的某些想法——尤其是「极限三角剖分」的概念——也许能为这些新的模函子提供一些有用的洞见。

\item 本文的故事最直接地适用于非共形的 $\N=2$ 理论。不过还有一些密切相关的共形理论（在 \cite{Gaiotto:2009we} 中进一步讨论），可以通过把 $C$ 上奇点处单值的本征值调到零来得到。在二维 $\N=(2,2)$ 理论的密切类比情形中，\cite{Cecotti:1993rm} 利用了关于大质量形变的信息来获取关于共形点的信息（例如其共形维数谱）。自然要问：类似的技巧在这里是否也行得通。

\item 最后，如上面所提，理解本文所构造的函数 $\CX_\gamma$ 在小 $R$ 时的解析结构将会非常有意思。

\end{itemize}
